\documentclass[11pt,a4paper]{article}
\usepackage[margin=0.85in,top=1in,bottom=1in]{geometry}
\usepackage{amsmath,amssymb,amsfonts}
\usepackage{graphicx}
\usepackage{float}
\usepackage{enumitem}
\usepackage{microtype}
\usepackage{caption}
\usepackage[hidelinks]{hyperref}
\usepackage[most,breakable]{tcolorbox}
\tcbuselibrary{breakable,skins}
\usepackage{fontspec}
\setmainfont{DejaVu Serif}
\setsansfont{DejaVu Sans}
\setmonofont{DejaVu Sans Mono}
\captionsetup{font=small,labelfont=bf,skip=4pt}
\setlist{leftmargin=1.5em,topsep=2pt}

%% Breakable problem card — prevents content cutoff across pages
\newtcolorbox{solcard}[3]{
  breakable, enhanced,
  colback=white, colframe=gray!50, arc=2pt,
  title={\textbf{#1}\hfill\textsf{\small #3\,pts}},
  fonttitle=\normalsize\sffamily\bfseries,
  before upper={\small\textit{#2}\par\smallskip\hrule height 0.4pt\smallskip},
  top=4pt, bottom=6pt, left=7pt, right=7pt,
  parbox=false,
}

\title{\textbf{Dynamics of quasiparticles --- Solution}}
\author{\normalsize Y23 (2023) — English translation}
\date{}
\begin{document}\maketitle\vspace{-0.5em}
\begin{solcard}{A1}{Let the dependence of electron energy on momentum have the form $$\varepsilon(p) = \varepsilon_0 (3 - \cos(ap_x) - \cos( ap_y) - \cos (ap_z)).$$ Find его acceleration at the moment времени, когда momentum equals $p_x = \dfrac{\pi}{2 a}$, $p_y = \dfrac{\pi}{3 a}$, $p_z = 0$, а компоненты действующей на electron силы $F_x$, $F_y$, $F_z$.}{0.70}

Let's express the speed of the electron in terms of its momentum: $$ v_x = \frac{\partial \varepsilon}{\partial p_x} = \varepsilon_0 a \sin (a p_x), $$similarly for остальных двух компонент $$ v_y = \varepsilon_0 a \sin(a p_y), \; v_z = \varepsilon_0 a \sin(a p_z). $$From the equation движения $\dfrac{d}{dt}\vec{p} = \vec{F}$. Let us express acceleration via/through производную momentumа: $$ a_x  = \dot{v}_x = \varepsilon_0 a^2 \cos (a p _x) \dot{p}_x = \varepsilon_0 a^2 \cos(a p_x) F_x. $$The formulas for the remaining components are similar.  For impulse values ​​from the condition

\par\smallskip\noindent\textbf{Answer:} \textbackslash{}textbf{Answer:} $$ a_x = 0, \; a_y = \frac{1}{2} \varepsilon_0 a^2 F_y, \; a_z = \varepsilon_0 a^2 F_z. $$
\end{solcard}

\begin{solcard}{A2}{Let an electron with $\varepsilon(\vec{p})$ from A1 move in a constant electric field of magnitude $E$ directed along the $x$ axis.  The initial coordinates and projections of the momentum are equal to zero. Electron charge $-q$. Find his law of motion $x(t)$.}{0.80}

The electron is acted upon by the force $F_x = - q E$, so the momentum $$ p_x = - q E t, $$а соответствующая velocity $$ v_x = -a \varepsilon_0 \sin(aqE t). $$Остальные проекции скорости equal нулю. Integrating, let us find зависимость координаты $x$ as a function of time

\par\smallskip\noindent\textbf{Answer:} \textbackslash{}textbf{Answer:} $$ x(t) = -\frac{\varepsilon_0}{q E}(1 - \cos(a q Et )) $$
\end{solcard}

\begin{solcard}{A3}{The concentration of free electrons in a certain substance is $n$, the dependence is $\varepsilon(p)$ as in A1. At the initial moment of time, a constant electric field $E$ is turned on, directed along the axis $x$. Find the dependence of the current density $j$ on time.}{0.50}

Current density is expressed in terms of electron speed and concentration $$ j_x = - n q v_x = n q a \varepsilon_0 \sin (a q E t). $$

\par\smallskip\noindent\textbf{Answer:} \textbackslash{}textbf{Answer:} $$ j_x = n q a \varepsilon_0 \sin (a q E t). $$
\end{solcard}

\begin{solcard}{B1}{Let the dependence of energy on momentum $\varepsilon(\vec{p})$ be arbitrary. Prove that the energy $\varepsilon$ and the projection of momentum onto the direction of the magnetic field are conserved.}{0.50}

The equation of motion is $$ \dot{\vec{p}} = - q \vec{v} \times \vec{B}. $$Производная энергии по времени $$ \frac{d \varepsilon} {dt} = \frac{\partial \varepsilon}{\partial p_x} \dot{p}_x + \frac{\partial \varepsilon}{\partial p_y} \dot{p}_y  +\frac{\partial \varepsilon}{\partial p_z} \dot{p}_z = v_x \dot{p}_x + v_y \dot{p}_y + v_z \dot {p}_z = \vec{v} \dot{\vec{p}}. $$Substituting expression for производной momentumа, we get $$ \dot{\varepsilon} = \vec{v} \cdot\left(- q \vec{v}  \times \vec{B}\right) = 0. $$Умножив equation движения скалярно на magnetic field we get $$ \frac{d}{dt} (\vec{B}{\vec{p}}) = \vec{B} \dot{\vec{p}} = - q \vec{B} (\vec{v} \times \vec{B}) = 0, $$that is, the projection of the impulse onto the magnetic field is constant.

\par\smallskip\noindent\textbf{Answer:} \textbackslash{}textbf{Answer:} $$ \dot{\varepsilon} = \vec{v} \cdot\left(- q \vec{v}  \times \vec{B}\right) = 0. $$
\end{solcard}

\begin{solcard}{B2}{Let the dependence of energy on momentum have the form $$\varepsilon = \frac{p_x^2}{2 m_x} + \frac{p_y^2}{2 m_y} .$$ Energy electronа equals $\varepsilon$. Find частоту движения $\omega$. Fromобtimesandте траекторию движения частицы на плоскости $xy$, indicate the characteristic dimensions and direction of movement. Place the origin of coordinates at the center of symmetry of the trajectory.}{1.10}

Let's write the equations of motion: \textbackslash{}begin{align*} m_x\textbackslash{}dot{v}_x \&= - q v_y B;\textbackslash{}\textbackslash{} m_y \textbackslash{}dot{v}_y \&= q v_x B. \textbackslash{}end{align*} Let's exclude $v_y$ from the first equation by differentiating it and expressing $\dot{v}_y$ from the second equation. We will receive $$ m_x \ddot{v}_x = - q \dot{v}_y B = - \frac{q^2 B^2}{m_y} v_x, $$that is $$ \ddot{v}_x = - \frac{q^2 B^2}{m_x m_y} v_x. $$Это equation гармонических колебаний  с частотой $$ \omega = \frac{qB}{\sqrt{m_x m_y}}. $$


The dependence of speed on time has the form $$ v_x = A \cos (\omega t + \varphi), $$where амплитуду $A$ можно найти, зная энергию electronа: $$ A = \sqrt{\frac{2\varepsilon}{m_x}}. $$Вторую компоненту скорости можно найти from the equation движения $$ v_y = - \frac{m_x}{q B} \dot{v}_x = \sqrt{\frac{2\varepsilon}{m_x}} \frac{\omega}{qB} \sin(\omega t + \varphi) = \sqrt{\frac{2 \varepsilon}{m_y}} \sin (\omega t + \varphi). $$Integrating, we get траекторию частицы \textbackslash{}begin{align*} x(t) \&= \textbackslash{}sqrt{\textbackslash{}frac{2 \textbackslash{}varepsilon}{m_x}} \textbackslash{}frac{1}{\textbackslash{}omega} \textbackslash{}sin(\textbackslash{}omega t + \textbackslash{}varphi) = \textbackslash{}frac{\textbackslash{}sqrt{2 m_y \textbackslash{}varepsilon}}{q B} \textbackslash{}sin (\textbackslash{}omega t + \textbackslash{}varphi);\textbackslash{}\textbackslash{} y(t ) \&= - \textbackslash{}frac{\textbackslash{}sqrt{2 m_x \textbackslash{}varepsilon}}{qB} \textbackslash{}cos (\textbackslash{}omega t + \textbackslash{}varphi). \textbackslash{}end{align*} Thus, траектория представляет собой эллипс с полуосями $$ A_x = \frac{\sqrt{2 m_y \varepsilon}}{q B},\; A_y = \frac{\sqrt{2 m_x \varepsilon}}{qB}. $$The movement occurs counterclockwise.

\par\smallskip\noindent\textbf{Answer:} \textbackslash{}textbf{Answer:} $$ \omega = \frac{qB}{\sqrt{m_x m_y}}, \; A_x = \frac{\sqrt{2 m_y \varepsilon}}{q B},\; A_y = \frac{\sqrt{2 m_x \varepsilon}}{qB}. $$
\begin{figure}[H]
  \centering
  \includegraphics[width=0.72\textwidth]{T2_S_EN_assets/fig1.jpg}
\end{figure}
\end{solcard}

\begin{solcard}{B3}{Find the hodograph of the electron's momentum vector (the trajectory along which the end of the momentum vector moves if it is plotted from one point). Draw it on the plane $p_x$, $p_y$. Please indicate typical dimensions.}{0.60}

Taking into account the electric field, the equations of motion are \textbackslash{}begin{align*} m_x\textbackslash{}dot{v}_x \&= -qE- q v_y B;\textbackslash{}\textbackslash{} m_y \textbackslash{}dot{v}_y \&= q v_x B. \textbackslash{}end{align*} Let us introduce a new variable $u_y$ using the relation $$ v_y = - \frac{E}{B} + u_y, $$then уравнения движения for/atмут такой же вид, как and в предыдущей части: \textbackslash{}begin{align*} m_x\textbackslash{}dot{v}_x \&= - q u_y B;\textbackslash{}\textbackslash{} m_y \textbackslash{}dot{u}_y \&=  q v_x B. \textbackslash{}end{align*} Начальные условия $$ v_x (0)  = 0 ,\; v_y(0) = 0, \; u_y(0) = \frac{E}{B}. $$Taking into account этого solution $$ v_x = -\sqrt{\frac{m_y}{m_x}}\frac{E}{B} \sin \omega t;\; v_y =- \frac{E}{B}(1- \cos \omega t). $$Then momentumы $$ p_x =- \sqrt{m_x m_y} \frac{E}{B} \sin \omega t, \; p_y = -m_y \frac{E}{B}(1 - \cos \omega t ). $$

\begin{figure}[H]
  \centering
  \includegraphics[width=0.72\textwidth]{T2_S_EN_assets/fig2.jpg}
\end{figure}
\end{solcard}

\begin{solcard}{B4}{Find the law of motion $x(t)$, $y(t)$.}{0.40}

Integrating the speeds from the previous point, we get

\par\smallskip\noindent\textbf{Answer:} \textbackslash{}textbf{Answer:} \textbackslash{}begin{align*} x(t) \&= -\textbackslash{}frac{m_y E}{qB^2}  (1 - \textbackslash{}cos \textbackslash{}omega t)\textbackslash{}\textbackslash{} y(t) \&=  - \textbackslash{}frac{E}{B} t + \textbackslash{}frac{\textbackslash{}sqrt{m_x m_y} E}{q B^2} \textbackslash{}sin \textbackslash{}omega t, \textbackslash{}; \textbackslash{}omega = \textbackslash{}frac{qB}{\textbackslash{}sqrt{m_x m_y}}. \textbackslash{}end{align*}
\end{solcard}

\begin{solcard}{B5}{If the energy is greater than the minimum value, the trajectories on the plane become open, and the particle can travel indefinitely in some direction. Find this minimum energy $\varepsilon_{min}$. At energy $\varepsilon > \varepsilon_{min}$, draw the trajectory and indicate the characteristic geometric dimensions.}{1.40}

Let us write down the equations of motion $$ \dot{\vec{p}} = - q \vec{v} \times \vec{B} = - q \dot{\vec{r}}\times \vec{B}. $$Integrating, we get $$ \Delta \vec{p} = - q \Delta \vec{r}  \times \vec{B}. $$Multiplying это equality vectorно на $\vec{B}$. $$ \Delta \vec{p} \times \vec{B} = - q(\Delta \vec{r} \times \vec{B}) \times \vec{B} = q B^2  \Delta \vec{r}. $$Thus, vector перемещения можно выtimesandть via/through change in momentum $$ \Delta \vec{r} = \frac{1}{qB^2} \Delta \vec{p} \times \vec{B}. $$ Consider сначала  в momentumном пространстве. From постоянства энергии it follows, что годограф momentumа  – surface $$ \frac{p^2_y}{2m_y} + \varepsilon_0 (1 - \cos ap_x) = \varepsilon = const. $$Equation этой поверхности можно представить в виде $$ p_y = \pm \sqrt{2 m_y (\varepsilon - \varepsilon _0 + \varepsilon _0  \cos a p_x)}. $$В том сrayае, if $\varepsilon \le 2 \varepsilon_0$, годограф представляет собой замкнутую кривую. Then траектория also будет замкнутой кривой. If however $\varepsilon > 2 \varepsilon_0$, годограф momentumа – незамкнутая кривая, направленная вдоль оси $x$. Then траектория движения частицы - кривая, направленная вдоль $y$, for/atчем зависимость $x(y)$ – periodandческая.  Let us find period этой функции and ширину полосы вдоль оси $x$, которую она занимает. $$ \Delta y  = \frac{1}{qB} \frac{2\pi}{a}, \; \Delta x = \frac{1}{qB} \Delta p_y = \frac{\sqrt{2 m_y}}{qB} (\sqrt{\varepsilon} - \sqrt{\varepsilon - 2 \varepsilon _0}). $$

\begin{figure}[H]
  \centering
  \includegraphics[width=0.72\textwidth]{T2_S_EN_assets/fig3.jpg}
\end{figure}
\end{solcard}

\begin{solcard}{C1}{First, we neglect the variance and consider the case $\alpha = 0$. At what angles $\theta_1$, $\theta_2$ is decay possible?}{0.50}

Let's write down the laws of conservation of energy and momentum: $$ u p = u q_1 + u q_2,\; \vec{p} = \vec{q}_1 + \vec{q}_2. $$Since length vectorа $\vec{p}$ is equal to the sum of the lengths of the remaining two vectors, all three vectors must be directed in the same direction.

\par\smallskip\noindent\textbf{Answer:} \textbackslash{}textbf{Answer:} $\theta_1 = \theta_2 = 0$
\end{solcard}

\begin{solcard}{C2}{Express the modulus of the second phonon $q_2$ in terms of $p$, $q_1$, $\theta_1$. Obtain the expression in the limit of small $\theta_1$ up to second order in $\theta_1$.}{0.50}

From the law of conservation of momentum $$ \vec{q}_2 = \vec{p} - \vec{q}_1. $$Hence $$ q_2^2 = p^2 + q_1^2 - 2 p q_1 \cos \theta_1 = p^2 + q_1^2 - 2 p q_1 (1 - \theta_1^2/2) =  (p - q_1)^2+ p q_1 \theta_1^2. $$Fromвлечем корень and expanding в ряд до требуемого порядка $$ q_2 = \sqrt{  (p - q_1)^2+ p q_1 \theta_1^2} \approx p - q_1 + \frac{p q_1}{2(p - q_1)} \theta_1^2 $$

\end{solcard}

\begin{solcard}{С3}{The angle dependence of the emitted phonon momentum at small angles has the form $$q_1(\theta_1) = A + B \theta_1.$$ Determine постоянные $A$, $B$. Express их via/through $p$, $u$, $\alpha$.}{1.20}

Let us write the law of conservation of energy $$ u p + \alpha p^3 = u q_1 + \alpha q_1^3 + u q_2 + \alpha q_2^3. $$Substituting expression for $q_2$: $$ u p + \alpha p^3 = u q_1 + \alpha q_1^3 + u \left(p - q_1 +  \frac{p q_1}{2(p - q_1)} \theta_1^2\right) + \alpha \left( p - q_1 + \frac{p q_1}{2(p - q_1)} \theta_1^2\right)^3. $$Раскроем скобки and оставим вклады порядка $\theta_1^2$: $$ \alpha p^3  = \alpha q^3 + \frac{u p q_1}{2 (p - q_1)} \theta_1^2 +\alpha (p - q_1)^3 + \frac{3 \alpha}{2} p q_1 (p - q_1)\theta_1^2. $$Перенесем terms, не зависящие от угла, налево: $$ \alpha (p^3 - q_1^3 - (p-q_1)^3) = \frac{u p q_1 \theta_1^2}{2(p - q_1)} \left( 1 + \frac{3 \alpha}{u } (p - q_1)^2\right).$$Вторым terms в скобках в правой части можно пренебречь в силу условия $\alpha p^3 \ll u p$, в левой части раскроем скобки and we get $$ 3 \alpha (p - q_1) p q_1 = \frac{u p q_1 \theta_1^2}{2(p - q_1)}. $$Hence $$ (p-q_1 )^2 = \frac{u }{6 \alpha} \theta_1^2, \; q_1 = p - \sqrt{\frac{u}{6 \alpha}} \theta_1. $$

\par\smallskip\noindent\textbf{Answer:} \textbackslash{}textbf{Answer:} $$ A = p, \; B = - \sqrt{\frac{u}{6 \alpha}}. $$
\end{solcard}

\begin{solcard}{C4}{At what sign of the constant $\alpha$ is decay possible?}{0.30}

The equation in the previous paragraph has a solution if $\alpha > 0$.

\par\smallskip\noindent\textbf{Answer:} \textbackslash{}textbf{Answer:} $\alpha > 0$
\end{solcard}

\begin{solcard}{C5}{Find the maximum angle $\theta_{max}$ at which a phonon can be emitted. Under what condition on $p$ can the phonon momentum be considered small?}{0.50}

The maximum angle is achieved at the minimum momentum of the emitted phonon $q_1 \to 0$. Then from the expression for $q_1$ we find $$ \theta_{\max} = \sqrt{\frac{6 \alpha}{u}} p. $$Вся problem решалась в for/atближении малых углов. Это approximation workет for/at условии $\theta_{max} \ll 1$, that is $$ p \ll \sqrt{\frac{u}{\alpha}}. $$

\par\smallskip\noindent\textbf{Answer:} \textbackslash{}textbf{Answer:} $$ \theta_{\max} = \sqrt{\frac{6 \alpha}{u}} p. $$
\end{solcard}

\begin{solcard}{C6}{Express the angle $\theta_2$ at which the second phonon is emitted in terms of $\theta_1$, $p$, $u$, $\alpha$. Consider the angles small, take into account the first-order contributions from $\theta_1$, $\theta_2$.}{1.00}

From the results $\mathrm{C3}$ it follows that the phonon momenta are expressed in terms of emission angles as $$ q_1 = p - \sqrt{\frac{u}{6 \alpha}} \theta_1, \; q_2 = p - \sqrt{\frac{u}{6 \alpha}} \theta_2. $$ Let us write проекцию закона сохранения momentumа на ось, перпендикулярную направлению движения исходного фонона: \textbackslash{}begin{align*} q_1 \textbackslash{}theta_1 \&= q_2 \textbackslash{}theta_2, \textbackslash{}\textbackslash{} p \textbackslash{}theta_1 -  \textbackslash{}sqrt{\textbackslash{}frac{u}{6 \textbackslash{}alpha}} \textbackslash{}theta_1^2  \&=  p \textbackslash{}theta_2 -  \textbackslash{}sqrt{\textbackslash{}frac{u}{6 \textbackslash{}alpha}} \textbackslash{}theta_2^2. \textbackslash{}end{align*} Перегруппируем terms: $$ p(\theta_1 - \theta_2) = \sqrt{\frac{u}{6 \alpha}} (\theta_1^2 -\theta_2^2), \; p = \sqrt{\frac{u}{6 \alpha}} (\theta_1 + \theta_2). $$

\par\smallskip\noindent\textbf{Answer:} \textbackslash{}textbf{Answer:} $$ \theta_2 = \sqrt{\frac{6 \alpha}{u}} p - \theta_1 $$
\end{solcard}

\end{document}